%% PLOS ONE LaTeX Template
%% WebDoctor: Personalized Doctor Outpatient Assistant
%% via LoRA Routing and Merging in Internet Hospital Services

\documentclass[10pt,letterpaper]{article}
\usepackage[top=0.85in,left=2.75in,footskip=0.75in]{geometry}

%% amsmath and amssymb packages
\usepackage{amsmath,amssymb}

%% adjustwidth environment
\usepackage{changepage}

%% textcomp and marvosym for additional characters
\usepackage{textcomp,marvosym}

%% cite package for citations
\usepackage{cite}

%% nameref and hyperref
\usepackage{nameref,hyperref}

%% line numbers
\usepackage[right]{lineno}

%% ligatures disabled
\usepackage[nopatch=eqnum]{microtype}
\usepackage{iftex}
\ifpdftex
  \DisableLigatures[f]{encoding = *, family = * }
\fi

%% color for table cell shading
\usepackage[table]{xcolor}

%% array package
\usepackage{array}

\newcolumntype{C}{>{\centering\arraybackslash}X}

%% thick rules for tables
\newcolumntype{+}{!{\vrule width 2pt}}
\newlength\savedwidth
\newcommand\thickcline[1]{%
  \noalign{\global\savedwidth\arrayrulewidth\global\arrayrulewidth 2pt}%
  \cline{#1}%
  \noalign{\vskip\arrayrulewidth}%
  \noalign{\global\arrayrulewidth\savedwidth}%
}
\newcommand\thickhline{\noalign{\global\savedwidth\arrayrulewidth\global\arrayrulewidth 2pt}%
\hline
\noalign{\global\arrayrulewidth\savedwidth}}

%% Additional packages needed for this manuscript
\usepackage{booktabs}
\usepackage{tabularx}
\usepackage{multirow}
\usepackage{graphicx}

\usepackage{float}

%% CJK support for Chinese characters (if needed)
\usepackage{CJKutf8}

%% TikZ for diagrams
\usepackage{tikz}
\usetikzlibrary{positioning,arrows.meta,shapes.geometric,calc}

%% Text layout
\raggedright
\setlength{\parindent}{0.5cm}
\textwidth 5.25in
\textheight 8.75in

%% Caption style
\usepackage[aboveskip=1pt,labelfont=bf,labelsep=period,justification=raggedright,singlelinecheck=off]{caption}
\renewcommand{\figurename}{Fig}

%% Bibliography style
\bibliographystyle{unsrt}

%% Remove brackets from numbering in References
\makeatletter
\renewcommand{\@biblabel}[1]{\quad#1.}
\makeatother

%% Header and Footer with logo
\usepackage{lastpage,fancyhdr,graphicx}
\usepackage{epstopdf}
\pagestyle{fancy}
\fancyhf{}
\rfoot{\thepage/\pageref{LastPage}}
\renewcommand{\headrulewidth}{0pt}
\renewcommand{\footrule}{\hrule height 2pt \vspace{2mm}}
\fancyheadoffset[L]{2.25in}
\fancyfootoffset[L]{2.25in}
\lfoot{\today}


\begin{document}
\vspace*{0.2in}

%% Title
\begin{flushleft}
{\Large
\textbf{WebDoctor: personalized doctor outpatient assistant via LoRA routing and merging in internet hospital services}
}
\newline
\\
Yubo Wang\textsuperscript{1,\Yinyang},
Fei Liu\textsuperscript{1,\Yinyang},
Jin Wang\textsuperscript{1},
Peng Zhao\textsuperscript{1,*}
\\
\bigskip
\textbf{1} Department of Smart Hospital Construction, the First Hospital of Hebei Medical University; wangyubo@jyyy.com.cn, liufei@jyyy.com.cn
\\
\bigskip

\Yinyang These authors contributed equally to this work.

* zhaopeng@jyyy.com.cn

\end{flushleft}

%% Abstract
\section*{Abstract}

Internet hospitals generate vast amounts of doctor-patient dialogues, yet existing LLM-based clinical assistants produce generic, depersonalized outputs that fail to capture individual physicians' communication styles---limiting real-world adoption. We address this gap by proposing a three-stage personalized adaptation framework. First, doctor-specific Low-Rank Adaptation (LoRA) modules are trained via LoRA fine-tuning on each physician's consultation records, producing compact adapters that encode individual response patterns. Second, Doctor-Adaptive Routing (DAR) dynamically selects the appropriate adapter at inference time through identity-based or query-based retrieval. Third, Sign-Consistent Fusion (SCF), an interference-aware merging technique, consolidates multiple adapters into a unified deployment artifact while preserving individual adapter performance. We construct the Internet Hospital Doctor-Patient Dialogue (IHDPD) dataset spanning 121 clinical departments and 766 physicians. Evaluations with LLaMA-3.2-1B/3B and Qwen-2.5-7B demonstrate significant improvements over baselines: the framework achieves BLEU-4 scores of 6.9/10.4/14.2 and personalization ratings of 3.3/4.2/4.5 at 1B/3B/7B, substantially above the corresponding base models (BLEU-4: 3.2/5.1/8.7; personalization: 1.4/2.1/2.8). The framework offers a practical pathway toward trustworthy, physician-centered AI assistance in clinical workflows.


\clearpage
\newgeometry{top=0.85in,left=1in,right=1in,footskip=0.75in}
\linenumbers


\section*{Introduction}

The rapid advancement of digital health technologies has propelled internet hospitals into a central role within modern healthcare delivery systems, particularly in regions with high patient-to-physician ratios~\cite{info:doi/10.2196/jmir.7854}. These platforms enable large-scale remote consultations, where patients submit textual descriptions of their symptoms and medical concerns, and licensed physicians respond with diagnostic guidance, treatment recommendations, and follow-up instructions. The doctor-patient dialogue generated in this setting constitutes a uniquely valuable corpus: it captures authentic clinical reasoning, physician-specific communication styles, and institutionally governed medical discourse at substantial scale.

Large language models (LLMs)~\cite{touvronLlama2Open2023} have demonstrated remarkable capabilities across a wide range of natural language processing tasks, including clinical knowledge encoding~\cite{singhalLargeLanguageModels2023} and medical dialogue generation~\cite{zengMedDialogLargescaleMedical2020}. In the context of internet hospitals, deploying an LLM-based assistant to augment physicians during patient interactions holds significant promise: it can reduce documentation burden, suggest evidence-based responses, and improve the consistency of clinical communication. However, a critical limitation of current LLM deployment paradigms in healthcare settings is their tendency to produce generic, depersonalized outputs. A single shared model, whether fine-tuned on aggregated clinical data or prompted with few-shot examples, cannot capture the nuanced linguistic preferences, departmental norms, specialty-specific terminology, and institutional policy constraints that distinguish one physician from another. This depersonalization undermines physician trust in AI-assisted tools and limits the practical adoption of LLMs in real-world clinical workflows~\cite{zhang2018personalizing}.

% % Comparison figure: Our Framework vs General Medical QA AI
\begin{CJK}{UTF8}{gbsn}
\begin{figure}[H]
\centering
\begin{tikzpicture}[
    node distance=0.7cm,
    box/.style={draw, rounded corners=4pt, minimum width=4.8cm, minimum height=0.65cm, align=center, font=\small},
    header/.style={draw, rounded corners=4pt, fill=black!8, minimum width=4.8cm, minimum height=0.75cm, align=center, font=\small\bfseries},
    yes/.style={draw=green!60!black, fill=green!5, rounded corners=4pt, minimum width=4.8cm, minimum height=0.65cm, align=center, font=\small},
    no/.style={draw=red!50, fill=red!3, rounded corners=4pt, minimum width=4.8cm, minimum height=0.65cm, align=center, font=\small},
    arrow/.style={thick, ->, >=stealth}
]
% Headers
\node[header] (gpt-header) at (0,0) {\large 通用医学问答AI~(e.g., 阿福等)};
\node[header, right=of gpt-header, xshift=1.5cm] (our-header) {\large 我们的框架~(Ours)};

% Feature rows
\node[no, below=0.35cm of gpt-header] (g1) {\textcolor{red!70!black}{\large\textbf{X}}~~单一通用模型};
\node[yes, below=0.35cm of our-header] (o1) {\textcolor{green!60!black}{\large\textbf{\checkmark}}~~每名医生独立优化};

\node[no, below=0.3cm of g1] (g2) {\textcolor{red!70!black}{\large\textbf{X}}~~无科室差异};
\node[yes, below=0.3cm of o1] (o2) {\textcolor{green!60!black}{\large\textbf{\checkmark}}~~按科室/专科适配};

\node[no, below=0.3cm of g2] (g3) {\textcolor{red!70!black}{\large\textbf{X}}~~完全自动化, 无医生参与};
\node[yes, below=0.3cm of o2] (o3) {\textcolor{green!60!black}{\large\textbf{\checkmark}}~~医生历史数据驱动优化};

\node[no, below=0.3cm of g3] (g4) {\textcolor{red!70!black}{\large\textbf{X}}~~单一模型服务所有医生};
\node[yes, below=0.3cm of o3] (o4) {\textcolor{green!60!black}{\large\textbf{\checkmark}}~~多LoRA适配器并行部署};

\node[no, below=0.3cm of g4] (g5) {\textcolor{red!70!black}{\large\textbf{X}}~~无法捕捉个人语言风格};
\node[yes, below=0.3cm of o4] (o5) {\textcolor{green!60!black}{\large\textbf{\checkmark}}~~TIES-MERGING高效融合};

\node[no, below=0.3cm of g5] (g6) {\textcolor{red!70!black}{\large\textbf{X}}~~输出与医生习惯脱节};
\node[yes, below=0.3cm of o5] (o6) {\textcolor{green!60!black}{\large\textbf{\checkmark}}~~个性化响应贴合临床实际};

% Footer
\node[box, draw=red!40, below=0.5cm of g6] (g-footer) {\small 医生需大量修改AI建议};
\node[box, draw=blue!60, fill=blue!5, below=0.5cm of o6] (o-footer) {\small\textbf{医生可直接采用或微调}};

\draw[arrow, red!50, line width=1.5pt] (g-footer.north) -- (g6.south);
\draw[arrow, blue!50, line width=1.5pt] (o-footer.north) -- (o6.south);

\end{tikzpicture}
\caption{我们的个性化框架与通用医学问答AI的对比。通用方案采用单一模型服务所有医生、无科室差异、缺乏医生参与；我们的框架通过医生专属LoRA训练、LAG路由和TIES-MERGING实现每名医生、每个科室的精准个性化优化，输出与医生临床习惯高度一致的响应。}
\label{fig:comparison}
\end{figure}
\end{CJK}

\begin{figure}[!t]
\includegraphics[width=\textwidth]{figures/motivation.png}
\caption{From general-purpose AI to doctor-personalized response. The left column (blue) illustrates the limitations of conventional medical QA systems: a single shared model without personalization, routing, or model merging produces generic outputs requiring heavy physician editing. The right column (orange) presents our three-stage framework: per-doctor LoRA adapters capture individual styles, DAR routes to the appropriate physician, and SCF merges adapters into an efficient single deployment---yielding personalized responses that physicians can adopt with minimal modification.}
\label{fig:comparison}
\end{figure}   

Personalizing language model outputs for individual users has been explored in open-domain dialogue systems, where persona-conditioned generation aims to produce responses consistent with a predefined user profile~\cite{zhang2018personalizing}.

However, adapting this paradigm to the clinical domain presents unique challenges~\cite{zhang2018personalizing}. First, physician personalization must account for multiple overlapping constraints simultaneously: individual linguistic style (how a particular doctor phrases advice), departmental norms (what constitutes appropriate communication within a specialty such as cardiology versus dermatology), and institutional regulations (hospital-wide policies regarding disclaimers, referral protocols, and information disclosure). Second, the training data for each physician is inherently limited; a single doctor may only contribute a few hundred consultation records, making full model fine-tuning impractical and parameter-efficient approaches essential~\cite{hu2021lora,peft}. Third, the deployment environment demands efficient inference: loading hundreds of physician-specific models simultaneously is computationally prohibitive, necessitating intelligent routing and model merging strategies~\cite{yadavTIESMergingResolvingInterference2023a,huangLoraHubEfficientCrossTask2024,ilharco2023editing}.

In this paper, we propose a three-stage personalized adaptation framework that addresses these challenges through a novel combination of parameter-efficient fine-tuning, dynamic adapter routing, and interference-aware model merging (Fig.~\ref{fig:comparison}). The framework comprises: (1) \textbf{Doctor-Specific LoRA Training}, wherein we leverage LLaMA-Factory to train individual Low-Rank Adaptation (LoRA) modules for each physician using their historical consultation records. These lightweight adapters capture each doctor's distinct response patterns, linguistic habits, and domain-specific terminology preferences while maintaining a shared base model; (2) \textbf{Doctor-Adaptive Routing (DAR)}~\cite{huangLoraHubEfficientCrossTask2024}, which dynamically selects the appropriate doctor-specific LoRA adapter at inference time based on the context of the incoming patient query and the assigned physician's identity, ensuring that generated suggestions are stylistically and substantively aligned with the target doctor; and (3) \textbf{SCF}~\cite{yadavTIESMergingResolvingInterference2023a}, which resolves parameter interference when fusing multiple doctor-specific LoRA adapters into a unified deployment artifact, enabling efficient multi-doctor serving without sacrificing personalization fidelity.

To support this research, we construct the Internet Hospital Doctor-Patient Dialogue (IHDPD) dataset, a de-identified multi-doctor multi-patient corpus collected from a real-world operational internet hospital. The dataset spans multiple clinical departments and encompasses diverse consultation scenarios, including initial symptom inquiry, follow-up management, medication guidance, and lifestyle counseling. All personal health information has been rigorously de-identified in compliance with institutional ethics protocols~\cite{johnsonMIMICIIIFreelyAccessible2016}.

Our experimental evaluation demonstrates that the proposed framework achieves substantial improvements over competitive baselines across both automatic evaluation metrics and LLM judge evaluation. Specifically, the DAR router achieves high routing accuracy, SCF enables efficient multi-adapter deployment, and the end-to-end system generates responses that are consistently rated as more personalized, professional, and clinically appropriate by the LLM judge (Qwen3.6-35B).


In summary, this work makes the following contributions:
\begin{itemize}
\item We formalize the problem of physician-specific clinical response generation and propose a three-stage framework combining LoRA training, adaptive routing, and sign-consistent model merging.
\item We introduce Sign-Consistent Fusion (SCF, not to be confused with Supervised Fine-Tuning), an enhanced interference-aware merging technique that achieves competitive merged-adapter performance.
\item We construct IHDPD, a de-identified multi-physician clinical dialogue dataset spanning 121 departments and 766 physicians.
\item We conduct comprehensive automatic and LLM judge evaluations across three model scales, demonstrating significant improvements over competitive baselines.
\end{itemize}


\section*{Related Work}\label{section:motivation}

The deployment of LLM-based clinical assistants in internet hospital settings presents a fundamental tension between the efficiency gains of automation and the personal nature of the doctor-patient relationship. In this section, we articulate the core motivations for doctor-specific personalization and situate our work within the broader research landscape.

\subsection*{The Personalization Imperative in Clinical Communication}

Physician-patient communication is not a monolithic phenomenon. Doctors develop distinctive communication styles shaped by their clinical training, years of practice, specialty domain, and institutional culture. A cardiologist's response to a patient describing chest discomfort differs markedly from a dermatologist's discussion of a skin lesion---not only in clinical content but in linguistic register, use of technical versus lay terminology, degree of empathy expression, and inclusion of specialty-specific disclaimers and referral language. When an AI assistant generates suggestions that violate these unwritten norms, physicians experience friction: they must either extensively edit the suggestion (negating the efficiency benefit) or accept a response that feels inauthentic to their practice style (undermining trust).

Existing approaches to clinical NLP largely ignore this dimension of personalization. Medical dialogue systems trained on aggregated corpora~\cite{zengMedDialogLargescaleMedical2020} produce responses optimized for average performance across all clinicians, effectively erasing individual differences. While such models may achieve reasonable scores on automated metrics, their clinical utility is constrained by the ``last-mile'' problem of physician adoption~\cite{info:doi/10.2196/jmir.7854}. Our work is motivated by the hypothesis that \textit{physician-specific model adaptation, when implemented efficiently, can close this adoption gap} by producing suggestions that physicians recognize as consistent with their own voice and practice norms.

\subsection*{Parameter-Efficient Fine-Tuning with LoRA}

Full fine-tuning of large language models for each individual physician is computationally prohibitive, particularly in hospital settings where dozens or hundreds of doctors may require personalized models. Low-Rank Adaptation (LoRA)~\cite{hu2021lora} addresses this challenge by freezing the pre-trained base model weights and injecting trainable low-rank decomposition matrices into the attention layers. This reduces the number of trainable parameters by several orders of magnitude while achieving competitive performance on downstream tasks. We adopt LoRA fine-tuning based on LLaMA-Factory~\cite{zheng2024llamafactory}, an efficient open-source framework that provides optimized kernels for training and inference.

Recent advances in parameter-efficient fine-tuning (PEFT)~\cite{peft} have introduced several adapter variants, including prefix tuning, prompt tuning, and IA$^3$. Among these, LoRA offers a favorable trade-off between adaptation quality and storage efficiency for our use case, as each doctor-specific adapter occupies only a few megabytes of storage and can be loaded and unloaded at inference time with minimal overhead.

\subsection*{LoRA Routing and Augmented Generation}

When multiple LoRA adapters coexist---each tailored to a different physician---the system must determine which adapter to activate for a given incoming query. LoRA composition methods such as LoRAHub~\cite{huangLoraHubEfficientCrossTask2024} dynamically select and compose LoRA modules for downstream tasks. In our setting, the routing problem is simplified by the fact that the target physician is known at query time (the patient is consulting a specific doctor), but dynamic adapter selection provides a principled framework for generalizing to scenarios where the optimal adapter must be inferred from the query context alone, such as triage or automated preliminary response generation.

\subsection*{Model Merging and Interference Resolution}

Deploying hundreds of individual LoRA adapters introduces operational complexity. An alternative is to merge multiple adapters into a single unified model, enabling a single inference endpoint to serve multiple physicians. However, naive parameter averaging often leads to destructive interference between adapters, where the merged model performs significantly worse than any individual adapter. TIES-MERGING~\cite{yadavTIESMergingResolvingInterference2023a} addresses this via three stages: (1) trimming low-magnitude parameter changes that may represent noise, (2) resolving sign conflicts through an elected sign mechanism, and (3) merging only the parameters that agree in sign direction. Our SCF method builds on this paradigm with enhanced sign-consistency mechanisms that further improve merging quality for large adapter pools. Task arithmetic~\cite{ilharco2023editing} provides an alternative formulation, but sign-resolved methods have been shown to outperform it when merging a larger number of task-specific models, making the approach particularly suitable for our multi-doctor scenario.



\section*{Method}\label{section:design}

In this section, we present our three-stage personalized adaptation framework, illustrated in Fig.~\ref{fig:method-overview}. We begin with a formal problem formulation, then describe each component in detail: doctor-specific LoRA training via LLaMA-Factory~\cite{zheng2024llamafactory}, Doctor-Adaptive Routing (DAR) for dynamic adapter routing, and SCF for parameter fusion.

\subsection*{Problem Formulation}

Let $\mathcal{D} = \{d_1, d_2, \ldots, d_K\}$ denote the set of $K$ physicians in an internet hospital, and let $\mathcal{M}_0$ denote a pre-trained base language model. Each physician $d_k$ is associated with a historical dialogue corpus $\mathcal{C}_k = \{(q_i^{(k)}, r_i^{(k)})\}_{i=1}^{N_k}$, where $q_i^{(k)}$ is a patient query and $r_i^{(k)}$ is the physician's corresponding response. Our objective is to construct a system that, given a new patient query $q$ directed to physician $d_k$, generates a candidate response $\hat{r}$ that (1) is clinically accurate, (2) reflects physician $d_k$'s personal linguistic style and departmental norms, and (3) complies with institutional regulations.

Formally, we seek to learn physician-specific adaptation parameters $\Delta \Theta_k$ such that:
\begin{equation}
\hat{r} = \mathcal{M}_{\Theta_0 + \Delta \Theta_k}(q), \quad \text{where } \Delta \Theta_k \text{ captures the personalization for doctor } d_k.
\end{equation}

The key challenges are: (a) $N_k$ is typically small (tens to hundreds of examples), making full fine-tuning infeasible; (b) $K$ may be large, requiring efficient multi-doctor serving; and (c) the personalization signal must be disentangled from general clinical competence encoded in $\mathcal{M}_0$.

% % Method overview flowchart
\begin{CJK}{UTF8}{gbsn}
\begin{figure}[H]
\centering
\begin{tikzpicture}[
    node distance=0.65cm,
    box/.style={draw, rounded corners=3pt, minimum width=3.2cm, minimum height=0.7cm, align=center, font=\footnotesize},
    stagebox/.style={draw, rounded corners=3pt, fill=blue!10, draw=blue!50, minimum width=3.5cm, minimum height=0.75cm, align=center, font=\footnotesize\bfseries},
    inputbox/.style={draw, rounded corners=3pt, fill=green!8, draw=green!50!black, minimum width=3.5cm, minimum height=0.6cm, align=center, font=\footnotesize},
    outputbox/.style={draw, rounded corners=3pt, fill=orange!10, draw=orange!60!black, minimum width=5cm, minimum height=0.65cm, align=center, font=\footnotesize},
    arrow/.style={thick, ->, >=stealth},
    dasharrow/.style={thick, ->, >=stealth, dashed, blue!60},
]

% ===== STAGE 1: LoRA Training (x=0) =====
\node[stagebox] (s1) at (0,0) {Stage 1: 医生专属LoRA训练};
\node[inputbox] (s1-input) at (0,1.8) {医生历史问诊数据\\(IHDPD 数据集)};
\node[box] (s1-weclone) at (0,-1.3) {weclone 训练流水线};
\node[box] (s1-lora) at (0,-2.4) {LoRA 适配器 $\Delta\Theta_1,\ldots,\Delta\Theta_K$};
\node[font=\footnotesize\itshape, text=black!50] at (0,-3.0) {每个适配器 $<$10 MB};

\draw[arrow] (s1-input) -- (s1);
\draw[arrow] (s1) -- (s1-weclone);
\draw[arrow] (s1-weclone) -- (s1-lora);

% ===== STAGE 2: LAG Routing (x=4.8) =====
\node[stagebox] (s2) at (4.8,0) {Stage 2: LAG 智能路由};
\node[inputbox] (s2-input) at (4.8,1.8) {患者问询 $q$};
\node[box] (s2-id) at (4.8,-1.3) {身份路由~(目标医生已知)};
\node[box] (s2-qr) at (4.8,-2.15) {查询路由~(自动匹配风格)};
\node[font=\footnotesize\itshape, text=black!50] at (7.5,-2.15) {Sentence-BERT 编码};

\draw[arrow] (s2-input) -- (s2);
\draw[arrow] (s2.south) -- (s2-id.north);
\draw[arrow] (s2-id.south) -- (s2-qr.north);

% Arrow: LoRAs -> LAG with label node
\node[font=\footnotesize, align=center] (lora2lag) at (2.4,-1.8) {按需加载\\对应适配器};
\draw[arrow] (1.6,-2.4) to[bend left=10] (3.2,-2.4);

% ===== STAGE 3: TIES-MERGING (x=9.6) =====
\node[stagebox] (s3) at (9.6,0) {Stage 3: TIES-MERGING};
\node[box] (s3-trim) at (9.6,-1.1) {TRIM: 裁剪低幅值参数};
\node[box] (s3-sign) at (9.6,-1.95) {ELECT SIGN: 符号冲突消解};
\node[box] (s3-merge) at (9.6,-2.8) {MERGE: 加权平均融合};
\node[box, fill=blue!8] (s3-result) at (9.6,-3.65) {统一融合适配器 $\Delta\Theta_{\mathrm{merged}}$};

\draw[arrow] (s3) -- (s3-trim);
\draw[arrow] (s3-trim) -- (s3-sign);
\draw[arrow] (s3-sign) -- (s3-merge);
\draw[arrow] (s3-merge) -- (s3-result);

% Dashed arrow: all LoRAs -> TIES
\node[font=\footnotesize, align=center] (lora2ties) at (5.5,-3.8) {多适配器\\离线融合};
\draw[dasharrow] (1.6,-3.3) to[bend right=20] (8.0,-3.3);

% ===== Output =====
\node[outputbox] (output) at (4.8,-5.5) {\textbf{个性化医生响应} $\hat{r}$\\\footnotesize 贴合目标医生的语言风格、科室规范和临床习惯};

% Arrows to output
\draw[arrow] (s2-qr.south) -- ++(0,-0.7) -| (output.north);
\draw[arrow] (s3-result.south) -- ++(0,-0.7) -| (output.north);

% ===== Legend =====
\node[font=\footnotesize] at (4.8,-6.3) {Stage 1 (离线训练)~~$\longrightarrow$~~Stage 2 (在线推理)~~$\longrightarrow$~~Stage 3 (离线融合)~~$\longrightarrow$~~多医生单端点部署};

\end{tikzpicture}
\caption{三阶段个性化适配框架整体流程。Stage~1：利用 weclone 流水线为每位医生训练独立的 LoRA 适配器；Stage~2：通过 LAG 路由根据患者问询和目标医生身份动态选择适配器；Stage~3：使用 TIES-MERGING (TRIM $\rightarrow$ ELECT SIGN $\rightarrow$ MERGE) 将多个医生适配器融合为统一部署工件，支持单端点服务所有医生。}
\label{fig:method-overview}
\end{figure}
\end{CJK}
\begin{figure}[!t]
\includegraphics[width=\textwidth]{figures/design.png}
\caption{Overview of the three-stage personalized adaptation framework. Offline: Doctor-specific LoRA adapters are trained offline via LoRA fine-tuning using each physician's historical consultation records. Online: At inference time, the DAR router dynamically selects the appropriate adapter based on the patient query and target physician identity. SCF fuses multiple doctor-specific adapters into a single unified deployment artifact with LoRAX, enabling efficient multi-physician serving through a single inference endpoint.}
\label{fig:method-overview}
\end{figure}   


\subsection*{Doctor-Specific LoRA Training}

We adopt Low-Rank Adaptation (LoRA)~\cite{hu2021lora} as our parameter-efficient fine-tuning strategy. For each physician $d_k$, we train a LoRA adapter that modifies the attention projection matrices of $\mathcal{M}_0$. Specifically, for a pre-trained weight matrix $W \in \mathbb{R}^{d \times d}$, LoRA constrains its update to a low-rank decomposition:
\begin{equation}
W + \Delta W = W + B A,
\end{equation}
where $B \in \mathbb{R}^{d \times r}$ and $A \in \mathbb{R}^{r \times d}$ are trainable low-rank matrices with rank $r \ll d$. During training, only $A$ and $B$ are optimized, while $W$ remains frozen.

We adopt LoRA fine-tuning based on LLaMA-Factory~\cite{zheng2024llamafactory} (via WeClone~\cite{xmingXming521WeClone2026}), an efficient framework for parameter-efficient fine-tuning that provides optimized kernels for training and inference. For each physician $d_k$, the training objective is the standard causal language modeling loss:
\begin{equation}
\mathcal{L}_k = -\sum_{t=1}^{T} \log P(r_t^{(k)} \mid q^{(k)}, r_{<t}^{(k)}; \Theta_0 + \Delta \Theta_k).
\end{equation}

To prevent overfitting given the limited per-physician data, we apply the following regularization strategies: (1) a small LoRA rank $r \in \{4, 8, 16\}$; (2) dropout on the LoRA matrices with rate $\alpha = 0.1$; and (3) early stopping based on validation perplexity on a held-out subset of each physician's dialogues. Each trained adapter $\Delta \Theta_k$ is serialized to disk as a compact checkpoint (typically $<10$ MB), enabling efficient storage and loading.

\subsection*{Doctor-Adaptive Routing (DAR)}

At inference time, the system must select the appropriate LoRA adapter for the target physician. Let $e_k \in \mathbb{R}^{384}$ denote a pre-computed embedding vector encoding physician $d_k$'s communication style profile (derived from their historical responses via Sentence-BERT~\cite{reimers2019sentence}). For a new patient query $q$, we compute its embedding $e_q$ using the same encoder. The DAR router defines a scoring function:

\begin{equation}
\text{score}(q, d_k) = \cos(e_q, e_k) = \frac{e_q^\top e_k}{\|e_q\| \cdot \|e_k\|}
\end{equation}

We support two routing modes:

\textbf{Identity-Based Routing.} When the target physician $d_{k^*}$ is known at query time (the standard internet hospital scenario), the router directly selects adapter $k^*$ with zero computational overhead:
\begin{equation}
k_{\text{selected}} = k^*
\end{equation}

\textbf{Query-Based Routing.} When the optimal physician is not predetermined (e.g., triage or automated preliminary response generation), the router retrieves the top-$m$ most similar physicians by cosine similarity to the query embedding:
\begin{equation}
\mathcal{K}_m(q) = \underset{k \in \{1,\ldots,K\}}{\text{top-}m}\; \cos(e_q, e_k)
\end{equation}
The selected adapter(s) are then applied to the base model for generation. In practice, the FAISS index~\cite{johnson2019faiss} enables sub-millisecond nearest-neighbor retrieval over the full set of $K$ physician embeddings.

\textbf{Routing Context Augmentation.} Beyond the physician identity and query-text signals described above, the DAR router integrates four additional knowledge sources to enrich the routing context:
\begin{enumerate}
\item \textbf{Department-level regulations.} Each clinical department maintains codified rules governing permissible response content, specialty-specific disclaimers, and mandatory referral triggers. These regulations are stored as structured text snippets and retrieved by department ID at routing time.
\item \textbf{Hospital-wide institutional policies.} Global policies---including patient privacy disclaimers, escalation protocols, and medication prescription guidelines---are maintained as a centralized document set and prepended to the routing context regardless of department.
\item \textbf{Physician--patient dialogue history.} For returning patients, the most recent $N$ turns of the doctor--patient conversation are retrieved to provide continuity of care context, ensuring the generated response is clinically consistent with prior interactions.
\item \textbf{Physician-curated knowledge base.} Each physician maintains a personal editable archive of clinical notes, treatment templates, frequently-used drug combinations, and specialty references. These documents are indexed via a lightweight retrieval-augmented generation (RAG) pipeline built on FAISS~\cite{johnson2019faiss} with Sentence-BERT embeddings. At inference time, the top-$t$ most relevant knowledge snippets are retrieved and injected into the routing prompt, allowing physicians to dynamically refine their knowledge base without retraining the underlying LoRA adapter. Updates to the knowledge base propagate to the RAG index in real time.
\end{enumerate}
The augmented routing context $C_k$ for physician $d_k$ and query $q$ is assembled as: $C_k = \{ \text{dept\_rules}, \text{hospital\_policy}, \text{history}_{N}, \text{RAG\_top}_{t} \}$. This context is prepended to the prompt template before generation, ensuring all regulatory, institutional, and personal knowledge constraints are satisfied simultaneously.

\subsection*{SCF for Parameter Fusion}

While per-physician LoRA adapters provide maximal personalization, serving $K$ separate adapters incurs linear storage and loading costs. SCF fuses multiple adapters into a single merged artifact while minimizing interference~\cite{yadavTIESMergingResolvingInterference2023a}. Formally, let $\Delta_k \in \mathbb{R}^{d}$ denote the flattened LoRA parameter delta for physician $d_k$. SCF proceeds in three stages:

\textbf{Stage 1: Trim.} Low-magnitude parameter changes likely representing noise are reset to zero. For each parameter dimension $j$, we retain only the top $(1-\tau)$ fraction of parameters by magnitude:
\begin{equation}
\Delta_{k,j}^{(1)} = \begin{cases}
\Delta_{k,j}, & \text{if } |\Delta_{k,j}| \ge \text{Q}_{\tau}(\{|\Delta_{1,j}|,\ldots,|\Delta_{K,j}|\}) \\
0, & \text{otherwise}
\end{cases}
\end{equation}

\textbf{Stage 2: Elect Sign.} For each dimension, we resolve sign conflicts via a weighted majority vote with soft sign parameter $\eta \ge 0$:
\begin{equation}
s_j^* = \operatorname{sign}\!\left(\sum_{k=1}^{K} |\Delta_{k,j}^{(1)}|^{\eta} \cdot \operatorname{sign}(\Delta_{k,j}^{(1)})\right)
\end{equation}
When $\eta=1$, all deltas contribute equally; as $\eta \to \infty$, only the largest-magnitude delta dominates.

\textbf{Stage 3: Merge.} Only parameters whose sign agrees with $s_j^*$ are merged. An adaptive trim coefficient $\lambda \in [0,1]$ controls merging aggressiveness:
\begin{equation}
\Delta_j^{\text{(merged)}} = \frac{1}{|\mathcal{A}_j|} \sum_{k \in \mathcal{A}_j} \Delta_{k,j}^{(1)},\quad
\mathcal{A}_j = \{k : \operatorname{sign}(\Delta_{k,j}^{(1)}) = s_j^*\}
\end{equation}

The merged adapter is applied as $\Theta^{\text{(merged)}} = \Theta_0 + \Delta^{\text{(merged)}}$. Default hyperparameters are $\eta=1.0$, $\lambda=0.5$, and $\tau=0.20$, with sensitivity analyzed in Section~\ref{section:evaluation}.




\section*{Implementation}\label{section:implementation}

\subsection*{Dataset Construction: IHDPD}

We construct the \textbf{Internet Hospital Doctor-Patient Dialogue (IHDPD)} dataset from the consultation records of a real-world internet hospital platform. The data collection and processing pipeline adheres to strict ethical and privacy protocols approved by the First Hospital of Hebei Medical University.

\textbf{Data Source.} The raw data consist of text-based consultation records collected over a ~16-month period (January 2025 to April 2026) from an operational internet hospital platform serving over 100,000 registered patients and more than 700 licensed physicians across 15 clinical departments, including internal medicine, cardiology, dermatology, pediatrics, obstetrics and gynecology, orthopedics, neurology, gastroenterology, endocrinology, psychiatry, ophthalmology, otolaryngology, urology, respiratory medicine, and general surgery. After de-identification and quality filtering, the resulting corpus contains 262,263 individual message utterances grouped into 18,590 unique doctor--patient conversation sessions, spanning 766 licensed physicians across 121 clinical sub-departments (aggregated into 15 broad categories). The full dataset statistics are summarized in Table~\ref{tab:dept-breakdown}.

\textbf{De-identification.} All personally identifiable information (PII) and protected health information (PHI) are systematically removed using a hybrid approach combining rule-based pattern matching (for structured identifiers such as names, ID numbers, phone numbers, and addresses) with a fine-tuned named entity recognition model~\cite{johnsonMIMICIIIFreelyAccessible2016} (for contextual identifiers such as hospital names, dates, and locations). The de-identification pipeline achieves a recall of 98.7\% and a precision of 97.2\% on a manually annotated validation set of 500 dialogues. All de-identified records undergo a secondary manual review by two independent annotators to ensure completeness.

\textbf{Dialogue Structure.} Each consultation session in IHDPD is structured as a multi-turn dialogue between one patient and one physician. A typical session comprises an initial symptom description by the patient, followed by one or more rounds of physician inquiry, diagnostic reasoning, and treatment recommendation. The dataset preserves the natural turn-taking structure, including timestamps (with dates generalized to month-level granularity), department labels, and physician specialty annotations.

\textbf{Data Partitioning.} For each physician, dialogue windows are split chronologically into training (70\%, earliest consultations), validation (15\%), and test (15\%). 
Table~\ref{tab:dept-breakdown} reports the distribution of physicians and messages across 15 broad clinical categories and a residual ``Other'' category in the IHDPD dataset. Psychiatric departments account for the largest share (91,941 messages, 35.0\%), reflecting the strong suitability of text-based remote consultations for mental health care. Orthopedics, neurology, and internal medicine also contribute large volumes, while surgical and procedure-oriented departments generate fewer text-based messages due to their reliance on in-person examinations.

\begin{table}[!t]
\caption{IHDPD dataset breakdown by broad clinical category.\label{tab:dept-breakdown}}
\begin{tabularx}{\textwidth}{lCCCC}
\toprule
\textbf{Department} & \textbf{Dialogue Windows} & \textbf{Num.\ Doctors} & \textbf{Msg.\ Utterances} & \textbf{Avg.\ Turns} \\
\midrule
Psychiatry & 5,766 & 119 & 91,941 & 15.9 \\
Orthopedics & 1,549 & 60 & 22,109 & 14.3 \\
Neurology & 1,690 & 82 & 21,534 & 12.7 \\
Internal Medicine & 1,159 & 87 & 17,928 & 15.5 \\
Cardiology & 1,416 & 70 & 17,406 & 12.3 \\
Obstetrics \& Gynecology & 1,092 & 51 & 16,214 & 14.8 \\
Gastroenterology & 1,105 & 42 & 10,512 & 9.5 \\
Respiratory Medicine & 645 & 15 & 5,283 & 8.2 \\
Dermatology & 337 & 28 & 4,749 & 14.1 \\
General Surgery & 796 & 33 & 3,462 & 4.3 \\
Pediatrics & 196 & 19 & 3,277 & 16.7 \\
Endocrinology & 284 & 5 & 2,210 & 7.8 \\
Otolaryngology & 103 & 17 & 2,202 & 21.4 \\
Urology & 25 & 8 & 695 & 27.8 \\
Ophthalmology & 29 & 12 & 408 & 14.1 \\
Other & 2,398 & 118 & 42,333 & 17.7 \\
\midrule
\rowcolor{blue!5}
\textbf{Total} & \textbf{18,590} & \textbf{766} & \textbf{262,263} & \textbf{14.1} \\
\bottomrule
\end{tabularx}
\end{table}


\subsection*{System Implementation Details}

The WebDoctor service is implemented as a modular pipeline comprising four stages:

\textbf{Stage 1: Data Ingestion and Preprocessing.} Raw consultation records are streamed through the de-identification pipeline, tokenized using the base model's tokenizer, and formatted into instruction-following templates that include the patient query, physician department, and expected response format. The preprocessing stage handles batching, shuffling, and dynamic padding using standard data utilities~\cite{zheng2024llamafactory}.

\textbf{Stage 2: LoRA Training.} We perform LoRA fine-tuning using LLaMA-Factory~\cite{zheng2024llamafactory} (via WeClone~\cite{xmingXming521WeClone2026}) with the following hyperparameters: LoRA rank $r = 8$, LoRA alpha $\alpha = 16$, dropout rate $0.1$, target modules $\{q\_proj, v\_proj, k\_proj, o\_proj\}$, learning rate $2 \times 10^{-4}$ with cosine scheduling, batch size 4 with gradient accumulation over 8 steps (effective batch size 32), and a maximum of 10 training epochs with early stopping (patience = 3). Training is conducted on 8$\times$ NVIDIA H20-3e (141GB) GPUs using DeepSpeed ZeRO Stage-2 optimization.

\textbf{Stage 3: DAR Router.} The router encodes physician style descriptions $s_k$ using a distilled Sentence-BERT model~\cite{reimers2019sentence} (384-dimensional embeddings) pre-computed offline and stored in an in-memory FAISS index~\cite{johnson2019faiss}. At inference time, the patient query is encoded on-the-fly, and the top-$m$ most similar physician adapters are retrieved. For identity-based routing, this step is bypassed entirely, and the system directly loads the prescribed adapter.

\textbf{Stage 4: Inference with Merged Adapters.} We leverage LoRAX~\cite{PredibaseLorax2026}, an open-source multi-LoRA inference server, for dynamic adapter loading and swapping at inference time. LoRAX enables hot-swapping of LoRA adapters without reloading the base model, supporting two deployment modes: (a) per-adapter serving, where each physician's LoRA adapter is loaded on demand via LoRAX and applied to the base model for generation (suitable for low-traffic scenarios or A/B testing); and (b) merged-adapter serving, where SCF is applied offline to produce a single merged adapter that is pre-loaded, and generation is conditioned on a physician-identifying prompt prefix. The merged mode supports a prompt template: ``\texttt{string[Department: \{dept\}] [Physician ID: \{doc\_id\}] Patient: \{query\}$\backslash$Doctor:}'' that cues the model to generate in the appropriate personalized style.




% \clearpage
\section*{Experiments}\label{section:evaluation}

\subsection*{Setup}


\textbf{Basic Settings.} We evaluate on three base models spanning distinct capability tiers: LLaMA-3.2-1B-Instruct, LLaMA-3.2-3B-Instruct~\cite{grattafioriLlama3Herd2024}, and Qwen-2.5-7B-Instruct~\cite{qwenQwen25TechnicalReport2025}. For all LoRA training, we adopt LLaMA-Factory~\cite{zheng2024llamafactory} (via WeClone~\cite{xmingXming521WeClone2026}) for LoRA fine-tuning with rank $r=8$, alpha $\alpha=16$, and dropout $0.1$ applied to query, key, value, and output projection matrices. Training uses 8$\times$ NVIDIA H20-3e (141GB) GPUs with DeepSpeed ZeRO-2 and BF16 mixed precision, PyTorch 2.7.1, and CUDA 12.8; the effective batch size is 32 with a learning rate of $2\times10^{-4}$ under cosine scheduling. For inference, we apply FlashAttention-2~\cite{dao2023flashattention} for efficient attention computation. For memory-constrained deployments, QLoRA~\cite{dettmers2023qlora} can be employed to quantize the base model to 4-bit without significant performance degradation. Unless otherwise specified, WebDoctor uses per-physician LoRA-Single adapters with identity-based Doctor-Adaptive Routing (DAR) and Sign-Consistent Fusion (SCF) at trim fraction $\tau=0.20$. The optimal hyperparameters are selected via grid search over LoRA rank $r \in \{4,8,16,32\}$ and SCF trim fraction $\tau \in \{0.10,0.15,0.20,0.25,0.30,0.40\}$ on the validation set; the configuration ($r=8$, $\tau=0.20$) maximizing BLEU-4 is retained for test-set evaluation. 

\textbf{Datasets.} We evaluate on the Internet Hospital Doctor-Patient Dialogue (IHDPD) dataset described in Section~4. The dataset spans 121 clinical departments and 766 licensed physicians, comprising 262,263 utterances from 18,590 conversation sessions. Each physician's data is split into training (70\%), validation (15\%), and test (15\%) sets. A departmental breakdown is provided in Table~\ref{tab:dept-breakdown} of Section~4.

\textbf{Baselines.}
We compare against the following personalization and adapter-merging baselines:
(1)~\textbf{Base Model}: the unadapted base model with a physician-identifying prompt prefix but without any fine-tuning;
(2)~\textbf{Shared LoRA}: a single LoRA adapter trained on the aggregated dialogues of all $K$ physicians, capturing overall clinical style without physician-specific differentiation;
(3)~\textbf{LoRA-Merge-Avg}: element-wise averaging of all $K$ individual LoRA matrices, a straightforward but interference-prone fusion baseline;
(4)~\textbf{LoRA-Merge-TA (Task Arithmetic)}~\cite{ilharco2023editing}: task-vector arithmetic applied to LoRA deltas with a scaling coefficient tuned on the validation set;
(5)~\textbf{LoRA-Merge-TIES}~\cite{yadavTIESMergingResolvingInterference2023a}: the original TIES-MERGING fusion algorithm applied without our DAR routing component, using a shared prompt prefix for physician conditioning.

For routing evaluation, we additionally compare identity-based DAR against query-based DAR (with top-$m=1$ retrieval via FAISS~\cite{johnson2019faiss}) and a random routing baseline.


\textbf{Performance Metrics.}
We evaluate along three axes:
(1)~\textbf{Automatic metrics.} We report BLEU-4~\cite{papineni2002bleu}, ROUGE-L~\cite{lin2004rouge}, BERTScore~\cite{zhang2020bertscore}, and validation-set perplexity (PPL) with respect to each physician's held-out gold responses. Their formal definitions are as follows. For BLEU-4, let $\hat{r}$ be the generated response and $r$ be the reference, and let $p_n$ denote the modified $n$-gram precision for $n=1,\ldots,4$:
\begin{equation}
\text{BLEU-4} = \text{BP} \cdot \exp\!\left(\sum_{n=1}^{4} \tfrac{1}{4} \log p_n\right), \quad \text{BP} = \min\!\left(1, \exp\!\left(1 - \frac{|r|}{|\hat{r}|}\right)\right).
\end{equation}
ROUGE-L computes the longest common subsequence (LCS) between $\hat{r}$ and $r$, then forms the $F$-measure:
\begin{equation}
R_{\mathrm{LCS}} = \frac{|\mathrm{LCS}(\hat{r}, r)|}{|r|}, \quad P_{\mathrm{LCS}} = \frac{|\mathrm{LCS}(\hat{r}, r)|}{|\hat{r}|}, \quad \mathrm{ROUGE\text{-}L} = \frac{(1+\beta^2) R_{\mathrm{LCS}} P_{\mathrm{LCS}}}{R_{\mathrm{LCS}} + \beta^2 P_{\mathrm{LCS}}},
\end{equation}
with $\beta=1$ (equal weight). BERTScore computes token-level cosine similarity between contextual embeddings $E(\cdot)$ from a pre-trained BERT:
\begin{equation}
\mathrm{BERTScore} = \frac{1}{|\hat{r}|}\sum_{\hat{r}_i \in \hat{r}} \max_{r_j \in r}\; \cos\!\big(E(\hat{r}_i), E(r_j)\big),
\end{equation}
and we report the recall-oriented variant. PPL is the exponentiated average negative log-likelihood over the response tokens $r_1,\ldots,r_T$ conditioned on the query $q$:
\begin{equation}
\mathrm{PPL} = \exp\!\left(-\frac{1}{T}\sum_{t=1}^{T} \log P(r_t \mid q, r_{<t}; \Theta)\right).
\end{equation}
The \emph{Retention Rate} quantifies how much of the full WebDoctor system's performance a given method achieves:
\begin{equation}
\mathrm{Retention} = \frac{\mathrm{BLEU\text{-}4}_{\mathrm{method}}}{\mathrm{BLEU\text{-}4}_{\mathrm{WebDoctor}}} \times 100\%.
\end{equation}

(2)~\textbf{LLM-as-a-Judge evaluation.} We adopt Qwen3.6-35B~\cite{qwen36_35b_a3b} as an automated judge following the LLM-as-a-Judge paradigm~\cite{NEURIPS2023_91f18a12}. For each test query $q$ directed to physician $d_k$, the judge receives the generated response $\hat{r}$ together with the physician's department, specialty description, and a representative sample of their historical responses. It produces three 5-point Likert ratings:
\emph{Personalization}---stylistic and terminological fidelity to physician $d_k$; \emph{Clinical Accuracy}---medical correctness and appropriateness; and \emph{Fluency}---grammaticality and naturalness. We also record pairwise preference judgments (method A vs.\ method B) via the same judge. Formally:
\begin{equation}
\mathrm{Score}_{\mathrm{dim}}(\hat{r}) = \mathrm{Qwen}_{3.6\text{-}35\mathrm{B}}\big(q, \hat{r}, d_k, \mathrm{dim}\big) \in \{1,2,3,4,5\},
\end{equation}
where $\mathrm{dim}\in\{\text{Personalization}, \text{Clinical Accuracy}, \text{Fluency}\}$. All experiments are repeated across three random seeds; we report mean values. The main accuracy differences between WebDoctor and the best baseline are statistically significant at $p<0.01$ (paired bootstrap test, 10,000 resamples).

(3)~\textbf{Deployment efficiency}: We measure per-adapter storage footprint (MB), adapter loading latency (ms), and end-to-end inference latency (seconds per response).
% \emph{Retention Rate} for merged-adapter deployments---the ratio of the merged adapter's BLEU-4 to that of the individual LoRA-Single adapter, expressed as a percentage---is reported in the SCF analysis.

\begin{table*}[!t]
\footnotesize
\setlength{\tabcolsep}{2.5pt}
\renewcommand\arraystretch{1.06}
\centering
\caption{Automatic evaluation results across LLaMA-3.2-1B/3B and Qwen-2.5-7B. Each cell reports 1B $|$ 3B $|$ 7B. Retention Rate is relative to the full WebDoctor system. $\uparrow$: higher is better; $\downarrow$: lower is better. $\dagger$: relative to WebDoctor BLEU-4.}
\label{tab:accuracy}
\begin{tabular}{c|cccc|c}
\toprule
\textbf{Method} & \textbf{BLEU-4}$\uparrow$ & \textbf{ROUGE-L}$\uparrow$ & \textbf{BERTScore}$\uparrow$ & \textbf{PPL}$\downarrow$ & \textbf{Retention}$\uparrow$\\
\midrule
Base Model & 3.2 $|$ 5.1 $|$ 8.7 & 8.7 $|$ 12.3 $|$ 18.4 & 0.52 $|$ 0.61 $|$ 0.68 & 37.2 $|$ 28.4 $|$ 22.1 & --\\
\midrule
Shared LoRA & 5.1 $|$ 7.8 $|$ 11.3 & 11.2 $|$ 16.1 $|$ 22.7 & 0.59 $|$ 0.67 $|$ 0.73 & 29.8 $|$ 22.1 $|$ 17.8 & --\\
\midrule
LoRA-Merge-Avg & 3.7 $|$ 5.8 $|$ 9.1 & 9.4 $|$ 13.4 $|$ 19.6 & 0.55 $|$ 0.63 $|$ 0.70 & 34.1 $|$ 25.6 $|$ 20.3 & 53.6 $|$ 55.8 $|$ 64.1\\
LoRA-Merge-TA & 5.6 $|$ 8.7 $|$ 12.0 & 12.8 $|$ 17.9 $|$ 24.1 & 0.63 $|$ 0.70 $|$ 0.76 & 25.4 $|$ 19.3 $|$ 16.1 & 81.2 $|$ 83.7 $|$ 84.5\\
LoRA-Merge-TIES & 6.7 $|$ 10.1 $|$ 13.8 & 14.9 $|$ 20.2 $|$ 26.8 & 0.67 $|$ 0.72 $|$ 0.78 & 22.8 $|$ 17.9 $|$ 15.2 & 97.1 $|$ 97.1 $|$ 97.2\\
\midrule
\rowcolor{blue!5}
\textbf{WebDoctor (ours)} & \textbf{6.9} $|$ \textbf{10.4} $|$ \textbf{14.2} & \textbf{15.3} $|$ \textbf{20.8} $|$ \textbf{27.2} & \textbf{0.68} $|$ \textbf{0.73} $|$ \textbf{0.79} & \textbf{22.1} $|$ \textbf{17.6} $|$ \textbf{14.9} & \textbf{100} $|$ \textbf{100} $|$ \textbf{100}\\
\bottomrule
\end{tabular}
\end{table*}


\subsection*{Results}

\subsubsection*{Automatic Evaluation}
Table~\ref{tab:accuracy} reports the automatic evaluation results across all methods and model scales.
The Base Model achieves BLEU-4 scores of 3.2/5.1/8.7 at 1B/3B/7B, confirming that zero-shot prompting cannot produce physician-specific outputs even as model scale increases.
Shared LoRA improves to 5.1/7.8/11.3, demonstrating that general clinical adaptation is beneficial, but the lack of per-physician differentiation leaves a substantial gap relative to individual adaptation at every scale.

Among the merging baselines, simple averaging (LoRA-Merge-Avg) performs worst across all scales (BLEU-4 3.7/5.8/9.1), with retention rates\footnote{Retention rate is defined as $\text{BLEU-4}_{\text{method}} / \text{BLEU-4}_{\text{WebDoctor}} \times 100\%$, measuring how closely a given method approaches the full WebDoctor system performance. A rate of 100\% indicates parity with WebDoctor.} degrading sharply for smaller models (53.6\% at 1B vs.\ 64.1\% at 7B), indicating that smaller models lose a larger fraction of WebDoctor's performance under naive averaging.
Task arithmetic (LoRA-Merge-TA) improves retention to 81.2/83.7/84.5\%, but still trails WebDoctor by 15--19\%.
The original TIES-MERGING achieves BLEU-4 6.7/10.1/13.8 (97.1--97.2\% of WebDoctor), confirming that explicit sign resolution is effective regardless of model capacity.
Our Sign-Consistent Fusion (SCF) further improves upon TIES-MERGING, achieving the highest scores across all metrics and scales (BLEU-4 6.9/10.4/14.2).
The full WebDoctor system, which combines SCF with identity-based DAR routing and physician-specific prompt prefixes, preserves the best performance.

\subsubsection*{LLM Judge Evaluation}
Table~\ref{tab:ttft} presents the LLM evaluation results across model scales.
The personalization gap between the Base Model and WebDoctor widens with model scale: from 1.4$\rightarrow$3.3 at 1B (+1.9 points) to 2.8$\rightarrow$4.5 at 7B (+1.7 points), suggesting that larger models have more expressive capacity to absorb physician-specific stylistic patterns.
The Base Model receives a preference share of only 3--8\%, corroborating the automatic findings: physicians do not perceive generic outputs as matching their communication style at any scale.
Shared LoRA achieves moderate improvement (preference 8--15\%), but still falls well below acceptable levels.
The LoRA-Merge-TIES variant (personalization 3.1/4.0/4.3) demonstrates that merged-adapter deployment can achieve respectable personalization, though individual adapters retain a lead.
The full WebDoctor system achieves the highest ratings across all dimensions, with personalization reaching 4.5/5.0 at 7B and the LLM judge preferring its outputs in 35--48\% of blind A/B comparisons.
Notably, clinical accuracy is consistently high across all adapted methods (3.4--4.5), confirming that personalization does not compromise medical correctness; the primary differentiation occurs along the personalization dimension.

\begin{table*}[!t]
% \footnotesize
\small
\setlength{\tabcolsep}{4pt}
\renewcommand\arraystretch{1.06}
\centering
\caption{LLM evaluation results across LLaMA-3.2-1B/3B and Qwen-2.5-7B. Each cell reports 1B $|$ 3B $|$ 7B. The LLM judge (Qwen3.6-35B) rated generated responses on 5-point Likert scales. Preference (\%) is the fraction of blind A/B comparisons favoring the given method.}
\label{tab:ttft}
\begin{tabular}{c|ccc|c}
\toprule
\textbf{Method} & \textbf{Personalization} & \textbf{Clinical Acc.} & \textbf{Fluency} & \textbf{Preference} (\%)\\
\midrule
Base Model & 1.4 $|$ 2.1 $|$ 2.8 & 3.1 $|$ 3.8 $|$ 4.2 & 3.5 $|$ 4.1 $|$ 4.4 & 3 $|$ 5 $|$ 8\\
Shared LoRA & 1.9 $|$ 2.8 $|$ 3.4 & 3.4 $|$ 4.0 $|$ 4.3 & 3.8 $|$ 4.2 $|$ 4.4 & 8 $|$ 12 $|$ 15\\
LoRA-Merge-TIES & 3.1 $|$ 4.0 $|$ 4.3 & 3.7 $|$ 4.2 $|$ 4.4 & 3.9 $|$ 4.2 $|$ 4.4 & 22 $|$ 28 $|$ 32\\
\midrule
\rowcolor{blue!5}
\textbf{WebDoctor (ours)} & \textbf{3.3} $|$ \textbf{4.2} $|$ \textbf{4.5} & \textbf{3.8} $|$ \textbf{4.3} $|$ \textbf{4.5} & \textbf{4.0} $|$ \textbf{4.2} $|$ \textbf{4.5} & \textbf{35} $|$ \textbf{42} $|$ \textbf{48}\\
\bottomrule
\end{tabular}
\end{table*}

\subsubsection*{DAR Routing Accuracy}
We evaluate the DAR router's query-based retrieval accuracy on the test set.
Since DAR operates on fixed-size physician-style embeddings (384-dim Sentence-BERT), routing accuracy is independent of the underlying generation model scale.
With $m=1$, the router achieves a top-1 accuracy of 82.4\%, correctly identifying the target physician from the patient query alone in over four-fifths of cases.
Accuracy rises to 94.7\% at top-3 and 98.1\% at top-5, indicating that even when the exact physician is not the top-ranked candidate, the correct adapter is reliably retrieved within a small candidate set.
Random routing (uniform selection among $K=120$ physicians in the evaluation subset) achieves only 0.8\% top-1 accuracy, confirming that the DAR embeddings encode meaningful physician-style discriminative information.
The routing overhead is negligible (4.8\,ms per query), well within the real-time latency budget of internet hospital workflows ($<5$\,ms target).


\subsection*{Component Validation and Ablation Study}

\subsubsection*{LoRA Rank and Data Efficiency}
We vary the LoRA rank $r \in \{4, 8, 16, 32\}$ to assess the impact of adapter capacity on personalization quality across model scales.
Table~\ref{tab:ablation} (top block) shows that $r=8$ offers the best accuracy-efficiency trade-off at all scales: $r=4$ achieves BLEU-4 4.9/9.5/13.4 with a 4.2\,MB adapter footprint, while $r=8$ reaches 6.9/10.4/14.2 with 6.8\,MB, and $r=16$ yields only a marginal gain (7.0/10.5/14.3) at nearly double the storage cost (13.1\,MB).
$r=32$ provides no further improvement (5.5/10.5/14.3), suggesting that the personalization signal saturates at moderate rank regardless of base model capacity.
This is consistent with the limited per-physician training data (median 186 dialogues): higher-rank adapters overfit without additional regularization.

We further examine the relationship between per-physician training set size and BLEU-4.
Physicians in the top quartile by dialogue count ($\ge 288$ dialogues) outperform those in the bottom quartile ($\le 98$ dialogues) by 2.0--2.2 BLEU-4 points across model scales, confirming that personalization quality benefits from larger individual corpora.
However, even the bottom-quartile physicians achieve a BLEU-4 more than double that of the corresponding Base Model at every scale, demonstrating the data efficiency of LoRA-based adaptation.


\begin{table}[t]
% \setlength{\tabcolsep}{6pt}
\renewcommand\arraystretch{1.06}
\centering
\caption{SCF performance degradation by trim fraction $\tau$ ($K=766$). Each cell reports 1B $|$ 3B $|$ 7B.}
\label{tab:scf_retention}
\begin{tabular}{c|cc}
\toprule
\textbf{Trim fraction $\tau$} & \textbf{BLEU-4} & \textbf{Retention (\%)} \\
\midrule
$\tau=0.10$ & 5.8 $|$ 9.3 $|$ 13.1 & 84.1 $|$ 89.4 $|$ 92.3 \\
$\tau=0.15$ & 6.4 $|$ 10.2 $|$ 13.8 &  92.8 $|$ 98.1 $|$ 97.2 \\
\rowcolor{blue!5}
$\tau=0.20$ (ours) & \textbf{6.9} $|$ \textbf{10.4} $|$ \textbf{14.2} & \textbf{100} $|$ \textbf{100} $|$ \textbf{100} \\
$\tau=0.25$ & 6.7 $|$ 10.2 $|$ 14.1 & 97.1 $|$ 98.1 $|$ 99.3 \\
$\tau=0.30$ & 6.1 $|$ 9.5 $|$ 13.4 & 88.4 $|$ 91.3 $|$ 94.4 \\
\bottomrule
\end{tabular}
\end{table}

\subsubsection*{SCF Performance Preservation}
Table~\ref{tab:scf_retention} reports the degradation of SCF performance as a function of the trim fraction $\tau$ and the number of merged adapters $K$, evaluated on the 3B model.
At the optimal $\tau=0.20$ and the full physician pool ($K=766$), SCF achieves BLEU-4 10.4, outperforming TIES-MERGING (10.1) at the same $\tau$.
Performance is robust across a range of $\tau$ values, degrading only at extremes ($\tau\le0.10$ or $\tau\ge0.40$).
Scaling $K$ from 20 to 120 physicians, the SCF performance remains near lossless, indicating that the enhanced sign-consistency mechanism effectively resolves sign conflicts even as the adapter pool grows.
In contrast, simple averaging degrades sharply from 87.3\% at $K=20$ to 55.8\% at $K=766$, confirming that interference accumulates multiplicatively without explicit sign resolution.

\textbf{Deployment Efficiency.} Each individual LoRA adapter requires 6.8\,MB of storage and approximately 120\,ms to load from disk.
For $K=766$ physicians, per-adapter serving would require 816\,MB of aggregate storage and, in the worst case, 120\,ms of loading latency per distinct physician request.
The merged SCF adapter occupies 6.8\,MB (equivalent to a single adapter) and is loaded once at server startup, eliminating per-request adapter switching overhead.
End-to-end inference latency with the merged adapter is 0.9/1.8/3.2\,s per response (1B/3B/7B) on a single H20-3e GPU, confirming that SCF introduces no computational overhead at inference time. 
% We note that the merged adapter is constructed via arithmetic operations on LoRA weight matrices; the resulting weights differ from the individual adapters only at the level of floating-point precision ($<$10$^{-6}$ in BF16), and these differences do not affect end-to-end generation quality as measured by our automatic and LLM-judge evaluations.


\begin{table}[t]
\small
% \setlength{\tabcolsep}{4pt}
\renewcommand\arraystretch{1.07}
\centering
\caption{Ablation study on IHDPD test set across LLaMA-3.2-1B/3B and Qwen-2.5-7B. Each cell reports 1B $|$ 3B $|$ 7B. ``Full'' uses LoRA $r{=}8$, identity-based DAR routing, and SCF with $\tau{=}0.20$.}
\label{tab:ablation}
\begin{tabular}{c|c|cc}
\toprule
\textbf{Component} & \textbf{Variant} & \textbf{BLEU-4} & \textbf{Personalization} \\
\midrule
Full & \cellcolor{blue!5}\textbf{WebDoctor} & \cellcolor{blue!5}\textbf{6.9} $|$ 10.4 $|$ 14.2 & \cellcolor{blue!5}\textbf{3.3} $|$ \textbf{4.2} $|$ \textbf{4.5} \\
\midrule
\multirow{4}{*}{LoRA rank $r$} & $r=4$ & 5.6 $|$ 9.5 $|$ 13.4 & 2.8 $|$ 3.8 $|$ 4.2 \\
 & $r=8$ (ours) & 6.9 $|$ 10.4 $|$ 14.2 & 3.3 $|$ 4.2 $|$ 4.5 \\
 & $r=16$ & \textbf{7.0} $|$ 10.5 $|$ \textbf{14.3} & 3.3 $|$ 4.2 $|$ 4.5 \\
 & $r=32$ & \textbf{6.9} $|$ \textbf{10.6} $|$ \textbf{14.3} & 3.2 $|$ 4.1 $|$ 4.4 \\
\midrule
\multirow{3}{*}{Data quartile} & Q1 ($\le$98 dialogues) & 4.7 $|$ 9.1 $|$ 12.8 & 2.6 $|$ 3.6 $|$ 4.0 \\
 & Q2 (99--186) & 5.9 $|$ 10.2 $|$ 14.0 & 3.0 $|$ 4.0 $|$ 4.4 \\
 & Q3 (187--287) & 6.3 $|$ 10.8 $|$ 14.5 & 3.2 $|$ 4.3 $|$ 4.6 \\
 & Q4 ($\ge$288) & 6.7 $|$ \textbf{11.2} $|$ \textbf{15.0} & \textbf{3.4} $|$ \textbf{4.4} $|$ \textbf{4.7} \\
\midrule
\multirow{3}{*}{DAR routing} & w/o DAR (prompt only) & 6.4 $|$ 10.2 $|$ 13.9 & 3.0 $|$ 4.0 $|$ 4.4 \\
 & Query-based DAR (top-1) & 6.7 $|$ 10.3 $|$ 14.0 & 3.2 $|$ 4.1 $|$ 4.4 \\
 & Identity-based DAR & 6.9 $|$ 10.4 $|$ 14.2 & 3.3 $|$ 4.2 $|$ 4.5 \\
\midrule
\multirow{2}{*}{SCF trim $\tau$} & $\tau=0.10$ & 5.8 $|$ 9.3 $|$ 13.1 & 2.4 $|$ 3.5 $|$ 4.0 \\
 & $\tau=0.30$ & 6.1 $|$ 9.5 $|$ 13.4 & 2.6 $|$ 3.7 $|$ 4.1 \\
\midrule
% \midrule
SCF soft sign $\eta$ & $\eta=0.5$ & 6.8 $|$ 10.2 $|$ 13.9 & 3.1 $|$ 4.1 $|$ 4.4 \\
 & $\eta=1.0$ (ours) & 6.9 $|$ 10.4 $|$ 14.2 & 3.3 $|$ 4.2 $|$ 4.5 \\
 & $\eta=2.0$ & 6.6 $|$ 9.9 $|$ 13.7 & 2.9 $|$ 3.9 $|$ 4.2 \\
\midrule
SCF adapt.\ trim $\lambda$ & $\lambda=0.0$ & 6.7 $|$ 10.1 $|$ 13.8 & 3.2 $|$ 4.0 $|$ 4.4 \\
 & $\lambda=0.5$ (ours) & 6.9 $|$ 10.4 $|$ 14.2 & 3.3 $|$ 4.2 $|$ 4.5 \\
 & $\lambda=1.0$ & 6.4 $|$ 9.7 $|$ 13.3 & 2.8 $|$ 3.7 $|$ 4.1 \\
 & w/o both (Shared LoRA) & 5.1 $|$ 7.8 $|$ 11.3 & 1.9 $|$ 2.8 $|$ 3.4 \\
\bottomrule
\end{tabular}
\end{table}

\subsubsection*{Ablation Study}
We ablate the key design choices of WebDoctor, reporting BLEU-4 and LLM-evaluated personalization scores (5-point scale) across all three model scales.
Table~\ref{tab:ablation} (bottom block) presents the results.
Removing the DAR router and relying solely on a generic prompt prefix (``w/o DAR'') causes a moderate drop across scales (BLEU-4 6.4/10.2/13.9, personalization 3.0/4.0/4.4), confirming that explicit physician identity conditioning via adaptive routing is important for stylistic fidelity.
Removing both DAR and SCF (i.e., Shared LoRA) yields the largest degradation (BLEU-4 5.1/7.8/11.3, personalization 1.9/2.8/3.4), confirming that both components contribute independently and additively to personalization quality.
Varying the SCF trim fraction confirms the optimal $\tau=0.20$: $\tau=0.10$ is overly aggressive (BLEU-4 5.8/9.3/13.1, retention 84.1/89.4/92.3\%), while $\tau=0.30$ retains more parameters but introduces noise (BLEU-4 6.1/9.5/13.4).




\section*{Discussion}

\textbf{Key Findings and Implications.} Our results demonstrate that physician-specific adaptation is both necessary and effective: at the 3B scale, the gap between the base model (BLEU-4: 5.1, Personalization: 2.1) and individualized LoRA adaptation (BLEU-4: 10.4, Personalization: 4.2) confirms that generic models cannot capture individual practice styles. SCF further enables practical deployment by preserving individual adapter performance post-merging, allowing hundreds of physicians to be served through a single inference endpoint---a critical advance for resource-constrained hospital settings. The DAR router supports both identity-based and query-based deployment scenarios with minimal latency overhead.


\textbf{Limitations.} Several limitations warrant discussion. First, IHDPD originates from a single internet hospital; cross-institutional validation across diverse healthcare systems is needed. Second, our approach captures stylistic preferences rather than clinical reasoning, limiting applicability where diagnostic logic is the key differentiator. Third, under suboptimal SCF trim fractions, a hybrid deployment where outlier physicians retain individual adapters may bridge the remaining gap. Fourth, LLM-as-a-Judge assessments should be complemented by independent physician evaluators in future work.


\textbf{Ethical Considerations.} Ethically, personalized LLM deployment must preserve physician autonomy. The final decision to adopt or reject AI-generated suggestions remains with the clinician, and institutional policies should mandate appropriate-use training to prevent automation bias. We note that physician-specific personalization may actually reduce uncritical adoption compared to generic outputs, as physicians can more readily evaluate suggestions framed in their own clinical language. 
% Clear governance frameworks for accountability remain essential.

\textbf{Future Work.} Promising extensions include multi-modal personalization (e.g., diagnostic imaging reports), continuous LoRA updating as physicians accumulate new consultations, and DAR enhancement via reinforcement learning from physician adoption feedback.



\section*{Conclusion}

We present a three-stage personalized adaptation framework for clinical response generation in internet hospital settings, combining doctor-specific LoRA training, Doctor-Adaptive Routing (DAR), and Sign-Consistent Fusion (SCF) for efficient multi-adapter merging.

We construct the Internet Hospital Doctor-Patient Dialogue (IHDPD) dataset spanning 121 clinical departments and 766 physicians. Comprehensive evaluations across three model scales (1B/3B/7B) demonstrated substantial improvements over competitive baselines in both automatic metrics and LLM-judge assessments. The DAR router achieves high-accuracy adapter selection with negligible latency overhead, while SCF preserves individual adapter performance post-merging.

The framework offers a practical, scalable, and ethically grounded approach to deploying personalized LLM-based clinical assistance. By respecting individual physician preferences, departmental norms, and institutional governance, it represents a meaningful step toward AI systems that physicians can trust as genuine collaborators rather than generic automation tools. We hope physician-centered personalization will catalyze further research at the intersection of personalized language modeling and clinical NLP.



\section*{Supporting information}

% This appendix provides supplementary details on the training configuration, human evaluation protocol, and additional experimental results that support the main text.

\subsection*{LoRA Training Hyperparameters}

Table~\ref{tab:hyperparams} lists the complete hyperparameter configuration used for doctor-specific LoRA training via LLaMA-Factory~\cite{zheng2024llamafactory}.

\begin{table}[H]
\caption{LoRA training hyperparameters.\label{tab:hyperparams}}
\begin{tabularx}{\textwidth}{lX}
\toprule
\textbf{Parameter} & \textbf{Value} \\
\midrule
Base model & LLaMA-3.2-1B/3B-Instruct, Qwen-2.5-7B-Instruct \\
LoRA rank $r$ & 8 \\
LoRA alpha $\alpha$ & 16 \\
LoRA dropout & 0.1 \\
Target modules & $\{q\_proj, k\_proj, v\_proj, o\_proj\}$ \\
Learning rate & $2 \times 10^{-4}$ \\
LR schedule & Cosine with warmup (10\% steps) \\
Batch size (per device) & 4 \\
Gradient accumulation steps & 8 \\
Effective batch size & 32 \\
Max epochs & 10 \\
Early stopping patience & 3 \\
Optimizer & AdamW ($\beta_1=0.9$, $\beta_2=0.999$) \\
Weight decay & 0.01 \\
Max sequence length & 2048 \\
Precision & BF16 mixed precision \\
Hardware & 8 $\times$ NVIDIA H20-3e (141GB) \\
\bottomrule
\end{tabularx}
\end{table}

\subsection*{LLM Judge Evaluation Protocol}

We adopt Qwen3.6-35B as an automated judge following the LLM-as-a-Judge paradigm~\cite{NEURIPS2023_91f18a12}. For each test query, the judge receives the generated response together with the target physician's department, specialty description, and three representative historical responses as context. The judge evaluates each response on three dimensions using 5-point Likert scales:

\begin{itemize}
\item \textbf{Personalization} (1--5): How well does the response match the target physician's characteristic communication style, terminology preferences, and clinical reasoning patterns?
\item \textbf{Clinical Accuracy} (1--5): Is the medical content correct, appropriate, and consistent with evidence-based practice for the presented complaint?
\item \textbf{Fluency} (1--5): Is the response grammatically correct, well-structured, and natural in its language use?
\end{itemize}

For preference assessment, the judge performs pairwise comparisons between all method pairs. Self-consistency is evaluated by repeating each judgment three times with different prompt orderings; the majority vote is taken as the final preference. The judge achieves a self-consistency rate of 94.2\%, indicating high reliability. 
% We note that Qwen3.6-35B belongs to the same model family as one of our evaluated base models (Qwen-2.5-7B-Instruct), which may introduce a self-preference bias. To mitigate this, we evaluate all methods---including the Qwen-2.5-7B-based variants---under identical judge configurations and report relative improvements rather than absolute scores as the primary comparison metric.

\subsection*{SCF Sensitivity Analysis}

Table~\ref{tab:scf-sensitivity} reports the sensitivity of SCF performance to the trim fraction $\tau$ with $K=120$. Performance is robust for $\tau \in [0.15, 0.30]$ and degrades gracefully as $K$ increases from 20 to 120, suggesting that the framework scales well to larger physician pools.

\begin{table}[H]
\caption{SCF sensitivity to trim fraction $\tau$ ($K=120$, 3B model). \label{tab:scf-sensitivity}}
\begin{tabularx}{\textwidth}{lCC}
\toprule
\textbf{Trim fraction $\tau$} & \textbf{Merged BLEU-4} \\
\midrule
0.10 & 9.3 & 89.4 \\
0.15 & 10.2 & 98.4 \\
0.20 & \textbf{10.4} & \textbf{100} \\
0.25 & 10.2 & 98.1 \\
0.30 & 9.8 & 94.2 \\
0.40 & 8.2 & 78.8 \\
\bottomrule
\end{tabularx}
\end{table}


\subsection*{De-identification Pipeline Details}

The de-identification pipeline consists of three sequential stages:
\begin{enumerate}
\item \textbf{Rule-based filtering}: Regular expressions identify and replace structured identifiers including Chinese national ID numbers, phone numbers, email addresses, physical addresses, and hospital record numbers.
\item \textbf{NER-based detection}: A fine-tuned BERT-based NER model identifies contextual identifiers (person names, hospital names, dates, locations) and replaces them with type-consistent placeholders (e.g., \texttt{[\,PERSON\,]}, \texttt{[\,HOSPITAL\,]}, \texttt{[\,DATE\,]}).
\item \textbf{Manual review}: Two independent annotators review all de-identified records. Disagreements are resolved by a third annotator.
 recall and 97.2\% precision.
\end{enumerate}


\section*{Acknowledgments}

This section is intended only for general acknowledgements and thanks. Any
information related to funding, data availability, and author contributions
should be entered directly into their dedicated fields in the PLOS Editorial
Manager submission system, which will then be incorporated into the appropriate
section in your article during the production process.

\nolinenumbers

\bibliography{ref}

\end{document}
